\documentclass[11pt]{article}
\usepackage[margin=25mm]{geometry}
\usepackage{amsmath,amssymb,amsthm}
\newtheorem{theorem}{Theorem}
\newtheorem{corollary}{Corollary}
\title{A lower bound for a cap at a singleton layer}
\author{Human Compression project}
\date{}
\begin{document}\maketitle

\begin{theorem}\label{cap}
Let $H$ be a finite connected simple graph of girth at least five.
Suppose that one vertex $r$ has degree two and every other vertex has
degree four. Then $|V(H)|\geq19$.
\end{theorem}
\begin{proof}
In a graph without triangles or quadrangles, the vertices reached by
nonbacktracking walks of lengths zero, one and two from a fixed vertex
are all distinct. Consequently, if a vertex has degree four and its
four neighbors also have degree four, its radius-two ball has
$1+4+4\cdot3=17$ vertices.

Write $a,b$ for the neighbors of $r$. The radius-two ball about $r$
has nine vertices: $r$, $a,b$, and two disjoint triples $P,Q$ of
children of $a,b$. In particular, a vertex of $P$ exists. It has
degree four and is not adjacent to $r$, so all its neighbors have
degree four. Its radius-two ball has seventeen vertices, and hence
$|V(H)|\geq17$.

There is therefore a vertex outside the nine-vertex radius-two ball
of $r$. Any such vertex has degree four, all its neighbors have degree
four, and its radius-two ball excludes $r$. Thus $|V(H)|\geq18$.

Suppose for a contradiction that equality holds. Let $C$ be the nine
vertices outside the radius-two ball of $r$. For every $x\in C$,
its seventeen-vertex radius-two ball consists of all vertices except
$r$. It therefore contains $a,b$. Since $x$ is outside the radius-two
ball of $r$, it is adjacent to neither $a$ nor $b$. A path of length
two from $a$ to $x$ must go through $P$, and one from $b$ to $x$
must go through $Q$. The vertex $x$ has at most one neighbor in each
triple: two neighbors in $P$ would form a quadrangle through $a$,
and similarly for $Q$. Thus $x$ has exactly one neighbor in each of
$P,Q$. This also proves that every vertex of $C$ has distance exactly
three from $r$; the layer sizes $(1,2,6,9)$ are a consequence, not an
additional assumption.

Label $P=\{p_0,p_1,p_2\}$ and $Q=\{q_0,q_1,q_2\}$, and label
$x\in C$ by $(i,j)$ when its parents are $p_i,q_j$. Two vertices of
$C$ cannot have the same label, since they and their two parents
would form a quadrangle. All nine labels therefore occur exactly once.
Every vertex of $C$ has its two remaining neighbors in $C$: it has
exactly two neighbors in $P\cup Q$, none in $\{r,a,b\}$, and degree
four. For completeness, the eighteen parent incidences from $C$
also consume all eighteen degrees left at $P\cup Q$ after their
edges to $a,b$, so there is no edge within $P\cup Q$.

Adjacent vertices of $C$ have different first coordinates and different
second coordinates, since a common parent would give a triangle.
Moreover the two $C$-neighbors of any vertex have different first
coordinates and different second coordinates: a common parent would
give a quadrangle through that vertex.

Choose an edge of $H[C]$, which exists because every vertex there
has degree two. Permuting the labels of $P$ and $Q$ separately, its
ends may be called $(0,0)$ and $(1,1)$. The other neighbor of $(0,0)$
must differ in both coordinates from each of these two labels, and
so is $(2,2)$. The same argument for $(1,1)$ makes its other neighbor
the same vertex $(2,2)$. These three vertices form a triangle, the
required contradiction.
\end{proof}

\begin{corollary}\label{layer}
Let $G$ be a finite connected simple four-regular graph of girth at
least five, with distance layers $V_0,\ldots,V_T$ from a fixed root.
If $V_i$ is a singleton and $0<i<T$, then
\[
 \sum_{h=0}^{i}|V_h|\geq19,
 \qquad
 \sum_{h=i}^{T}|V_h|\geq19.
\]
\end{corollary}
\begin{proof}
There are no edges between layers differing by more than one. The
boundary of any vertex set in a four-regular graph is even, since
$|\partial S|=4|S|-2e(G[S])$. The singleton has a positive even
number of edges into each neighboring side, because both proper
distance prefixes have nonempty boundaries. Its degree is four,
so these numbers are both two.

Each of the subgraphs induced by $V_0\cup\cdots\cup V_i$ and by
$V_i\cup\cdots\cup V_T$ is therefore a cap of the kind in
Theorem~\ref{cap}: the singleton has internal degree two, and all
other vertices retain degree four. Both subgraphs are connected.
For the prefix this follows from root-directed shortest paths; for
the suffix the same paths reach the unique vertex of $V_i$ without
leaving the suffix. Girth at least five is inherited by induced
subgraphs. Apply the theorem to each cap.
\end{proof}

\paragraph{Application to the retained order-94 profile.}
The exact scalar profile in
\texttt{singleton\_layer\_section\_application.json} has $T=22$,
$n_{19}=1$, and terminal sizes
\[
 (n_{19},n_{20},n_{21},n_{22})=(1,2,6,9).
\]
Their sum is eighteen, violating Corollary~\ref{layer}. This profile
cannot be the distance-layer sequence of the stipulated graph. Its
previously verified satisfaction of the older scalar inequalities
and the between-singletons section bound remains correct.

\paragraph{Scope and attribution.}
The radius-two count and the triangle/quadrangle parent restrictions
are elementary, and no historical novelty is claimed for them.
The cap argument is a further necessary restriction for the existing
F009 target. It excludes this one order-94 scalar profile, not every
order-94 graph or every scalar profile. No existence claim for a cap
of order nineteen or twenty follows. This proof uses no enumeration,
solver, external classification, compilation, or Lean verification.
It changes no accepted enumeration prefix, source disposition, or
frozen computational request. The three repeated interior sections
may satisfy their separately accepted bounds while this final cap
fails; those sections are not invalidated by the obstruction.

\end{document}
